\documentclass{beamer}

\usepackage{amsmath,amsfonts,amssymb,bm}
\usepackage{physics}
\usepackage{braket}

\usetheme{Madrid}

\title{Can Transformers Be Used to Solve Partial Differential Equations?}
%\subtitle{A Physics-Oriented Perspective}
\author{Morten Hjorth-Jensen}
\date{FYS5429/9429 Spring 2026}

\begin{document}

\frame{\titlepage}

%================================================
\section{Main Idea}
%================================================

\begin{frame}{Short Answer}
Yes, transformers can be used to solve partial differential equations (PDEs), but not usually in the same direct sense as a classical numerical solver.

\medskip

Instead, they typically learn an \emph{operator map}:
\[
\mathcal{L}[u](x,t) = f(x,t)
\qquad \Longrightarrow \qquad
u = \mathcal{G}(f),
\]
where \(\mathcal{G}=\mathcal{L}^{-1}\) is the solution operator.

\medskip

The transformer learns an approximation
\[
\mathcal{G}_\theta : f \mapsto u.
\]
\end{frame}

%------------------------------------------------

\begin{frame}{Physical Interpretation}
A PDE defines a map from:
\begin{itemize}
\item coefficients,
\item forcing terms,
\item initial and boundary conditions
\end{itemize}
to a solution field.

\medskip

Thus the learning problem is:
\[
u(x,t) \approx \mathrm{Transformer}_\theta(\text{input fields}).
\]

From a PDE point of view, this is best seen as \emph{learning the inverse operator} associated with the PDE.
\end{frame}

%================================================
\section{Transformers and PDE Operators}
%================================================

\begin{frame}{Why Transformers Are Relevant}
Transformers are sequence-to-sequence models based on attention.

\medskip

For PDEs, one reinterprets:
\begin{itemize}
\item input tokens \(\rightarrow\) spatial or spacetime grid points,
\item hidden states \(\rightarrow\) local field values,
\item attention weights \(\rightarrow\) couplings between points.
\end{itemize}

This makes transformers natural candidates for learning nonlocal mappings.
\end{frame}

%------------------------------------------------

\begin{frame}{Attention as a Learned Kernel}
A transformer computes outputs of the general form
\[
u(x) \approx \sum_{x'} \alpha(x,x')\, v(x'),
\]
where \(\alpha(x,x')\) is a learned attention weight.

\medskip

Compare this with an integral representation of a PDE solution:
\[
u(x) = \int G(x,x') f(x')\,dx'.
\]

This suggests the correspondence
\[
\boxed{\text{attention} \;\sim\; \text{learned Green's function or kernel}.}
\]
\end{frame}

%------------------------------------------------

\begin{frame}{Operator Learning Viewpoint}
Many PDEs define solution operators
\[
\mathcal{G} : a(x), f(x), u_0(x) \mapsto u(x,t),
\]
where \(a(x)\) may be coefficients, \(f(x)\) forcing, and \(u_0(x)\) initial data.

\medskip

A transformer can be trained to approximate this map:
\[
u_\theta = \mathcal{G}_\theta(a,f,u_0).
\]

So one is not solving the PDE from scratch each time, but learning the family of solutions.
\end{frame}

%================================================
\section{Ways to Use Transformers for PDEs}
%================================================

\begin{frame}{Three Main Strategies}
There are several ways transformers can be used for PDEs:

\begin{enumerate}
\item \textbf{Purely data-driven operator learning}
\item \textbf{Physics-informed training}
\item \textbf{Hybrid numerical-ML methods}
\end{enumerate}
\end{frame}

%------------------------------------------------

\begin{frame}{1. Data-Driven Operator Learning}
Train on supervised pairs
\[
(\text{input fields}, \text{solution fields}).
\]

Examples:
\begin{itemize}
\item input: coefficient field + forcing
\item output: PDE solution
\end{itemize}

Then the model approximates
\[
u \approx \mathcal{G}_\theta(\text{input}).
\]

This is similar in spirit to neural operators.
\end{frame}

%------------------------------------------------

\begin{frame}{2. Physics-Informed Transformers}
Instead of relying only on solution data, one can enforce the PDE residual directly.

\medskip

If \(u_\theta\) is the transformer prediction, define a loss such as
\[
\mathcal{L}_{\mathrm{PDE}} =
\| \mathcal{L}[u_\theta] - f \|^2.
\]

Possible additions:
\begin{itemize}
\item boundary-condition losses,
\item initial-condition losses,
\item conservation-law penalties.
\end{itemize}

This is analogous to physics-informed neural networks (PINNs), but with transformer architectures.
\end{frame}

%------------------------------------------------

\begin{frame}{3. Hybrid Approaches}
Transformers can also be used together with classical PDE solvers.

Examples:
\begin{itemize}
\item accelerate coarse-to-fine corrections,
\item learn closures for multiscale models,
\item post-process or denoise numerical solutions,
\item emulate expensive substeps in time evolution.
\end{itemize}

In this setting the transformer acts more like a surrogate or correction model than a full solver.
\end{frame}

%================================================
\section{Advantages}
%================================================

\begin{frame}{Why Use Transformers?}
Transformers may be attractive for PDEs because they can:
\begin{itemize}
\item capture long-range and nonlocal interactions,
\item handle parametric families of PDEs,
\item learn reusable surrogate operators,
\item work on irregular dependencies more flexibly than some convolutional models.
\end{itemize}
\end{frame}

%------------------------------------------------

\begin{frame}{High-Dimensional and Parametric PDEs}
Transformers are particularly interesting when solving a PDE repeatedly for many parameter values:
\[
\mathcal{L}_\mu u = f,
\]
where \(\mu\) labels material parameters, geometry, or external fields.

\medskip

Instead of solving each instance separately, one trains a model to learn
\[
\mu \mapsto u_\mu.
\]

This can be valuable for uncertainty quantification, inverse problems, and control.
\end{frame}

%------------------------------------------------

\begin{frame}{Nonlocal Structure}
Many important physical systems have nonlocal couplings.

Examples:
\begin{itemize}
\item integral equations,
\item long-range interactions,
\item multiscale transport,
\item effective coarse-grained dynamics.
\end{itemize}

Attention is naturally suited to such settings because it directly couples distant points in the input.
\end{frame}

%================================================
\section{Limitations}
%================================================

\begin{frame}{Important Limitations}
Transformers do not automatically provide exact PDE solutions.

Potential issues include:
\begin{itemize}
\item lack of rigorous error guarantees,
\item costly training,
\item uncertain extrapolation outside the training set,
\item possible violation of symmetries or conservation laws,
\item reduced reliability for stiff or highly oscillatory systems.
\end{itemize}
\end{frame}

%------------------------------------------------

\begin{frame}{Generalization and Physical Structure}
A transformer may interpolate well within the training distribution while failing outside it.

\medskip

For physics applications, one must think carefully about:
\begin{itemize}
\item symmetry constraints,
\item conservation laws,
\item stability,
\item boundary conditions,
\item physically meaningful inductive biases.
\end{itemize}
\end{frame}

%================================================
\section{Comparison with Other ML Methods}
%================================================

\begin{frame}{Comparison with PINNs and Neural Operators}
Three important families are:

\begin{itemize}
\item \textbf{PINNs}: enforce the PDE residual directly
\item \textbf{Fourier Neural Operators (FNOs)}: learn in spectral space
\item \textbf{Transformers}: use attention to learn flexible nonlocal mappings
\end{itemize}

Each method has strengths and weaknesses depending on the structure of the PDE.
\end{frame}

%------------------------------------------------

\begin{frame}{Transformers vs Fourier Neural Operators}
Fourier neural operators are often especially strong when:
\begin{itemize}
\item the domain is regular,
\item spectral structure is important,
\item convolution-like global operators are natural.
\end{itemize}

Transformers may be preferable when:
\begin{itemize}
\item geometry is irregular,
\item interactions are complex or heterogeneous,
\item adaptive nonlocal couplings matter.
\end{itemize}
\end{frame}

%================================================
\section{Connections to Green's functions}
%================================================

\begin{frame}{Connection to Green's Functions}
A PDE solution can often be written formally as
\[
u = \mathcal{L}^{-1} f.
\]

In many cases this becomes
\[
u(x) = \int G(x,x') f(x')\, dx',
\]
where \(G(x,x')\) is a Green's function.

\medskip

This is why transformers can be interpreted as learning the kernel \(G(x,x')\) from data.
\end{frame}

%================================================
\section{Concrete Example}
%================================================

\begin{frame}{Example: Poisson Equation}
Consider the Poisson equation
\[
-\nabla^2 u(x) = f(x).
\]

Formally, the solution is
\[
u(x) = \int G(x,x') f(x')\, dx'.
\]

A transformer-based solver would:
\begin{itemize}
\item take discretized \(f(x)\) as input,
\item output a predicted field \(u(x)\),
\item learn the effective kernel linking the two.
\end{itemize}
\end{frame}

%------------------------------------------------

\begin{frame}{What the Transformer is Really Learning}
The key point is:

\[
\boxed{
\text{The transformer is not solving the PDE in the classical deterministic sense.}
}
\]

Instead, it is learning the map
\[
f \mapsto u,
\]
or more generally
\[
(a,f,u_0) \mapsto u.
\]

From a PDE viewpoint, this means learning an approximate Green's function or response operator.
\end{frame}

%================================================
\section{Summary}
%================================================

\begin{frame}{Summary}
Transformers can be used for PDEs.

\medskip

Their main role is to:
\begin{itemize}
\item learn PDE solution operators,
\item approximate Green's functions or inverse operators,
\item serve as fast surrogates for repeated PDE solves.
\end{itemize}

Conceptually:
\[
\text{transformer attention}
\;\longleftrightarrow\;
\text{learned kernel}
\;\longleftrightarrow\;
\text{approximate Green's function}.
\]
\end{frame}

%------------------------------------------------

\begin{frame}{Final Perspective}
A concise way to summarize the situation is:
\[
\boxed{
\text{Transformers do not usually solve PDEs exactly; they learn the mapping defined by the PDE.}
}
\]

Or, in more PDE-like language:
\[
\boxed{
\text{They learn an approximate inverse operator or Green's function.}
}
\]
\end{frame}

\end{document}
